\chapter{Build and compile order}
\label{ch:buildorder}

This chapter states the order in which the deliverable is compiled, why that
order is fixed, and what to do when a module is added. It is a short chapter,
but it documents a real failure mode: a design that compiles today because the
simulator happened to tolerate the order it was given, and stops compiling on
the next tool, the next flow, or after the next file is added.

\section{Why the order is fixed at all}

Verilog--2005 does not require a module to be compiled after the modules it
instantiates. Instances are resolved at \emph{elaboration}, not at compile time,
so a filelist in arbitrary order normally works. That is exactly what makes the
issue worth documenting rather than ignoring: the flow appears healthy while
resting on a property of one particular tool.

Two things genuinely are resolved at compile time, in file order.

\begin{enumerate}[leftmargin=*,itemsep=3pt]
\item \textbf{Macros and includes.} A \verb|`define| is visible only from the
point it is defined onwards within the compilation unit. Every module in this
design that uses a shared encoding includes \file{defines.vh} itself, under an
include guard, and the filelists carry the \texttt{+incdir+} that makes it
findable. That is what keeps this class of failure away --- not the ordering by
itself, but the ordering plus the discipline that no file depends on a macro
another file happened to define earlier.

\item \textbf{Name collisions.} A module defined twice, or a parameter name
clashing across files, is reported against whichever copy the compiler saw
first. In an arbitrary order the error points at an arbitrary file.
\end{enumerate}

Three further reasons are about engineering rather than about the language.

\begin{itemize}[leftmargin=*,itemsep=3pt]
\item A bottom-up order makes a genuine error legible. If a module is missing
from the list, the failure is reported at the level that needed it. In an
arbitrary order the same mistake surfaces as one unresolved instance at the top
of the hierarchy, with no indication of which level dropped it.
\item Several downstream tools --- lint, synthesis front-ends, formal
front-ends --- \emph{do} require definition before use. Keeping the simulation
filelist in that order means the same list can be handed to them unchanged.
\item The level a module sits at is a statement about the design. Reading
\file{rtl.f} tells a reviewer, before opening any source file, that
\reg{cpu\_top} is the only module with a dependency on anything other than
leaves.
\end{itemize}

\begin{rtlnote}
The rule is stated once and enforced in one place. \file{debug/sim/rtl.f} is the
single source of truth for the RTL source set and its order, and both the
ModelSim flow (\texttt{vlog -f rtl.f}) and the Verilator flow
(\texttt{verilator -f rtl.f}) read the same file. A module is added to the
project in exactly one location.
\end{rtlnote}

\section{Dependency levels}

The rule is: \emph{a module may only instantiate modules from a strictly lower
level}. Within a level, order is irrelevant and is kept alphabetical.

\begin{table}[H]
\centering
\caption{RTL compile order by dependency level (\file{debug/sim/rtl.f})}
\label{tab:rtllevels}
\begin{tabular}{@{}clL{0.44\linewidth}@{}}
\toprule
\textbf{Level} & \textbf{Modules} & \textbf{Instantiates} \\
\midrule
\multirow{14}{*}{L0}
 & \file{alu.v}              & --- \\
 & \file{axi\_lite\_slave.v} & --- \\
 & \file{branch\_predictor.v}& --- \\
 & \file{branch\_unit.v}     & --- \\
 & \file{control.v}          & --- \\
 & \file{csr\_file.v}        & --- \\
 & \file{decode.v}           & --- \\
 & \file{exception\_unit.v}  & --- \\
 & \file{fetch\_unit.v}      & --- \\
 & \file{hazard\_unit.v}     & --- \\
 & \file{imm\_gen.v}         & --- \\
 & \file{lsu.v}              & --- \\
 & \file{regfile.v}          & --- \\
 & \file{writeback\_mux.v}   & --- \\
\midrule
\multirow{3}{*}{L1}
 & \file{alu\_top.v} & \file{alu} \\
 & \file{mtimer.v}   & \file{axi\_lite\_slave} \\
 & \file{pic.v}      & \file{axi\_lite\_slave} \\
\midrule
L2 & \file{cpu\_top.v} & \file{alu\_top} (L1) and the twelve L0 modules other
than \file{alu} and \file{axi\_lite\_slave} \\
\bottomrule
\end{tabular}
\end{table}

\begin{figure}[H]
\centering
\begin{tikzpicture}[node distance=6mm]
  \node[blk,minimum width=112mm] (l0)
    {\textbf{L0} --- leaf modules, no submodule instances\\[2pt]
     \footnotesize alu \quad axi\_lite\_slave \quad branch\_predictor \quad branch\_unit \quad control\\
     \footnotesize csr\_file \quad decode \quad exception\_unit \quad fetch\_unit \quad hazard\_unit\\
     \footnotesize imm\_gen \quad lsu \quad regfile \quad writeback\_mux};
  \node[blkf,below=10mm of l0,minimum width=112mm] (l1)
    {\textbf{L1} --- depend on L0 only\\[2pt]
     \footnotesize alu\_top $\rightarrow$ alu \qquad mtimer $\rightarrow$ axi\_lite\_slave
     \qquad pic $\rightarrow$ axi\_lite\_slave};
  \node[blk,below=10mm of l1,minimum width=112mm] (l2)
    {\textbf{L2} --- system top\\[2pt]
     \footnotesize cpu\_top $\rightarrow$ alu\_top (L1) + twelve L0 modules};
  \draw[flow] (l0) -- (l1);
  \draw[flow] (l1) -- (l2);
  \node[right=3mm of l0.east,font=\scriptsize,align=left,rotate=0] {};
  \node[font=\scriptsize,left=1mm of $(l0.south)!0.5!(l1.north)$] {compiled first};
  \node[font=\scriptsize,left=1mm of $(l1.south)!0.5!(l2.north)$] {then};
\end{tikzpicture}
\caption{Compile order. Each level is compiled completely before the next one
begins, so every module is defined before anything instantiates it.}
\label{fig:compileorder}
\end{figure}

\section{Testbench and verification collateral}

The same rule extends to the verification code, split across three filelists so
that each flow takes only what it can compile.

\begin{table}[H]
\centering
\caption{Verification collateral, in compile order}
\label{tab:tblevels}
\begin{tabular}{@{}lL{0.30\linewidth}L{0.36\linewidth}@{}}
\toprule
\textbf{Filelist} & \textbf{Contents, in order} & \textbf{Used by} \\
\midrule
\file{rtl.f} & the design, L0 $\rightarrow$ L1 $\rightarrow$ L2 &
  both flows \\
\file{tb\_cpu.f} & \file{ck\_rst\_tb}, \file{axi\_lite\_mem\_model},
  \file{axi\_lite\_monitor}, \file{axi\_lite\_dec2}, then \file{tb\_cpu\_axi}
  which instantiates all four & both flows \\
\file{tb\_block.f} & the seven block-level benches, each of which
  instantiates exactly one RTL block & ModelSim only \\
(in \file{compile.do}) & \file{axi\_lite\_arb2}, then \file{tb\_dual\_core} &
  ModelSim only \\
(in \file{run\_verilator.sh}) & the SVA and coverage modules, then the
  \sig{bind} file last & Verilator only \\
\bottomrule
\end{tabular}
\end{table}

Two orderings inside that table are load-bearing rather than cosmetic. The
arbiter must precede the dual-core bench that instantiates it, for the same
legibility reason as the RTL. The \sig{bind} file must be compiled
\emph{last} of all, because a bind statement names both the target module and
the checker module and needs both to exist.

\section{Adding a module}

\begin{enumerate}[leftmargin=*,itemsep=2pt]
\item Place it at the lowest level that is still strictly above everything it
instantiates.
\item Insert it alphabetically within that level in \file{rtl.f}, with the
one-line comment naming what it instantiates.
\item If it introduces a new level, add the level header rather than widening
an existing one.
\end{enumerate}

Nothing else changes: \file{compile.do}, \file{regress.do} and
\file{run\_verilator.sh} all read the filelists, so neither flow is edited.

\section{Build entry points}

\begin{table}[H]
\centering
\caption{Build and run commands}
\label{tab:buildcmds}
\begin{tabular}{@{}L{0.42\linewidth}L{0.44\linewidth}@{}}
\toprule
\textbf{Command} & \textbf{What it does} \\
\midrule
\texttt{vsim -c -do "do compile.do; quit -f"} & compile only, in the documented
  order; must report \texttt{Errors: 0, Warnings: 0} \\
\texttt{vsim -c -do "do sim.do; quit -f"} & one run of the system bench \\
\texttt{vsim -c -do "do regress.do; quit -f"} & the full twelve-run regression
  from a single compile \\
\texttt{make modelsim} & the same regression, with the test programs
  reassembled first \\
\texttt{make test} & the Verilator assertion and coverage flow \\
\bottomrule
\end{tabular}
\end{table}

A clean compile is part of the pass criterion, not a courtesy: every source file
carries its own \verb|`timescale| so the result never depends on filelist order,
and every comparison, index and assignment is width-explicit rather than relying
on implicit extension or truncation.
